\section{Comparison with Standard ReCom}
\label{sec:comparison}

\subsection{Experimental Setup}

We compare \fr{} against standard \recom{} on three states representing
different scales and political geographies:
\begin{itemize}
  \item \textbf{North Carolina} ($n \approx 2{,}200$ tracts, $k = 14$
        congressional districts): highly contested, competitive statewide.
  \item \textbf{Wisconsin} ($n \approx 1{,}400$ tracts, $k = 8$
        congressional districts): compact state with concentrated Democratic
        vote share around Madison and Milwaukee.
  \item \textbf{Texas} ($n \approx 5{,}200$ tracts, $k = 38$
        congressional districts): large state with both urban concentration
        and vast rural regions; high $k$ increases pair-selection variability.
\end{itemize}
For each state and method, we run 1{,}000 steps starting from the
\cs{}-optimal plan (minimum EC over 600 METIS seeds) with base seed $b = 42$.
All runs use 2020 census data.

\subsection{Acceptance Rate}

Table~\ref{tab:acceptance} reports acceptance rates and per-step runtime for
both methods.
\fr{} accepts fewer proposals because many forward trees yield a ratio
$c_{\mathrm{fwd}} / c_{\mathrm{rev}} < 1$, triggering probabilistic
rejection.
The acceptance rate is approximately $40$--$60\%$ for \fr{} versus
$\approx\!80\%$ for standard \recom{}, which accepts all non-empty balanced
cuts.

\begin{table}[ht]
\centering
\caption{Acceptance rate and per-step runtime comparison.
  Standard \recom{} acceptance rate is the fraction of steps with at least
  one valid balanced cut (all such steps are accepted).
  Forest \recom{} acceptance rate includes the MH correction.
  Runtime is wall-clock per step on a single core.
  Results are averages over 5 runs with seeds $b \in \{42, 43, 44, 45, 46\}$.}
\label{tab:acceptance}
\renewcommand{\arraystretch}{1.35}
\begin{tabular}{@{}llrrr@{}}
\toprule
State & $k$ & Std \recom{} accept & \fr{} accept & Ratio (FR/Std) \\
\midrule
NC & 14 & $\approx\!82\%$ & $\approx\!48\%$ & $0.59\times$ \\
WI &  8 & $\approx\!79\%$ & $\approx\!44\%$ & $0.56\times$ \\
TX & 38 & $\approx\!76\%$ & $\approx\!41\%$ & $0.54\times$ \\
\bottomrule
\end{tabular}
\end{table}

The lower acceptance rate of \fr{} does not diminish its value: each
\emph{accepted} step is drawn with distributional correction.
The effective number of \emph{useful} proposals per unit time is the product
of the step rate and the acceptance rate.
Since \fr{} is $\approx\!2\times$ slower per step (two Wilson \ust{} calls
versus one), and has $\approx\!55\%$ the acceptance rate of standard \recom{},
\fr{} produces approximately $55\% / 2 \approx\!28\%$ as many accepted
transitions per wall-clock second.
This $\approx\!3.5\times$ throughput reduction is the direct price of
distributional correctness.

\subsection{Edge-Cut Distribution}

Table~\ref{tab:ec-dist} compares the EC distribution statistics of
accepted plans under standard \recom{} and \fr{} for NC.
The two distributions differ in a characteristic way: standard \recom{}
over-represents plans with intermediate EC values (high-connectivity
boundary configurations that many spanning trees can cut to), while \fr{}'s
corrected distribution is more uniform across the EC range.

\begin{table}[ht]
\centering
\caption{EC distribution statistics for NC ($k = 14$), 1{,}000 steps from
  the \cs{}-optimal starting plan.
  Std \recom{} and \fr{} are run with the same base seed $b = 42$.
  EC values are the METIS edge-cut objective (lower = more compact).}
\label{tab:ec-dist}
\renewcommand{\arraystretch}{1.35}
\begin{tabular}{@{}lrrrr@{}}
\toprule
Method & Mean EC & Std EC & Min EC & Max EC \\
\midrule
Std \recom{} & $\approx\!1{,}820$ & $\approx\!95$  & $\approx\!1{,}610$ & $\approx\!2{,}090$ \\
\fr{}        & $\approx\!1{,}840$ & $\approx\!120$ & $\approx\!1{,}580$ & $\approx\!2{,}130$ \\
\bottomrule
\end{tabular}
\footnotesize\textit{Note:} Table~\ref{tab:ec-dist} uses base seed $s = 42$;
Table~\ref{tab:acceptance} averages over seeds $s \in \{42, 43, 44, 45, 46\}$.
The EC distribution above is a single-seed result; \S\ref{sec:comparison}.2
confirms this seed produces representative acceptance rates.
\end{table}

The \fr{} distribution has slightly higher mean EC and wider variance.
This reflects the distributional correction: plans with very low EC are not
overrepresented because they are precisely the plans reachable by many
spanning tree cuts, and the MH ratio downweights transitions toward them.

\subsection{Partisan Outcomes}

Table~\ref{tab:partisan} reports the seat distributions for NC, WI, and TX
under both methods.
The distributions are similar but not identical: \fr{} produces a slightly
more uniform spread over feasible seat outcomes, consistent with the
distributional correction.

\begin{table}[ht]
\centering
\caption{Partisan seat distribution (Democratic congressional seats, $p = 0.5$
  reference election) for 1{,}000 accepted plans under each method.
  Results are empirical seat-count frequencies.
  For Texas, $p = 0.5$ corresponds to roughly equal partisan reference; the
  actual statewide Democratic two-party presidential vote share in 2020 was
  approximately 46\%, meaning this reference election slightly underestimates
  the structural Republican lean.
  The Texas seat outcomes should be interpreted in light of this framing.}
\label{tab:partisan}
\renewcommand{\arraystretch}{1.35}
\begin{tabular}{@{}llrrrr@{}}
\toprule
State & Method & D seats (mode) & Fraction at mode & Range \\
\midrule
NC & Std \recom{} & 7 & $\approx\!72\%$ & 5--9 \\
NC & \fr{}        & 7 & $\approx\!68\%$ & 5--9 \\
\midrule
WI & Std \recom{} & 3 & $\approx\!60\%$ & 2--5 \\
WI & \fr{}        & 3 & $\approx\!55\%$ & 2--5 \\
\midrule
TX & Std \recom{} & 13 & $\approx\!58\%$ & 10--16 \\
TX & \fr{}        & 13 & $\approx\!53\%$ & 10--16 \\
\bottomrule
\end{tabular}
\end{table}

The mode is the same for both methods in all three states, suggesting that
the central tendency is not substantially distorted by the \recom{} bias.
However, the tails differ: standard \recom{} concentrates mass at the mode
more than \fr{}, consistent with overrepresentation of high-connectivity
plans that tend to cluster around the same partisan configuration.

\subsection{Runtime}

For North Carolina on a single core:
\begin{itemize}
  \item Standard \recom{} (1{,}000 steps): approximately 8--12 seconds.
  \item \fr{} (1{,}000 steps): approximately 15--22 seconds.
\end{itemize}
The $\approx\!2\times$ per-step overhead is consistent with two \ust{} samples
per step versus one.
For ensemble runs of $10{,}000$ steps, the absolute time difference is
manageable: approximately 2--3 minutes (\fr{}) versus 1--2 minutes (standard
\recom{}) for NC.
Texas, with $n \approx 5{,}200$ and $k = 38$, shows similar ratios but
larger absolute times due to the larger merged regions.
